% Mobile and Embedded Computing, Lab 6. Compile: latexmk -xelatex lab06.tex
\documentclass[10pt,aspectratio=169,t]{beamer}
\usetheme[lab]{upbminimal}

\course{Mobile and Embedded Computing}
\decklabel{Lab 6}
\title{gRPC: a Go server and a Dart client}
\subtitle{One unary call, one stream, a deadline, and a cancel}
\author{Dinu-Ștefan Rusu}
\institute{FILS, Universitatea Politehnica București}
\date{Week 12}

\begin{document}

\titleframe

\begin{frame}[eyebrow=Today]{What this lab is for}
  \vskip 2mm
  \begin{grid}[2]
    \cell[\faServer]{Call a Go server once}{Generate a Dart client from a .proto file, call a unary RPC with a deadline, and read back a DeadlineExceeded status when the server is slow.}
    \cell[\faSatelliteDish]{Watch a stream live}{Subscribe to a server-streaming RPC, show one reading a second on screen, and cancel it cleanly with a button.}
  \end{grid}
  \vskip 6mm
  \taskmeta{Time}{2 hours}{Submit}{Push to your team repo, branch \code{lab-6}, before next lab}
\end{frame}

\begin{frame}[fragile,eyebrow=Setup]{Install the toolchain}
  \vskip 1mm
  Install these once per machine, then branch the team repository for this lab.
\begin{shell}
brew install protobuf   # or: apt install protobuf-compiler

go install google.golang.org/protobuf/cmd/protoc-gen-go@latest
go install google.golang.org/grpc/cmd/protoc-gen-go-grpc@latest

dart pub global activate protoc_plugin
\end{shell}
  \begin{hint}[Add them to PATH]
    The Go plugins land in \code{GOPATH/bin}, the Dart one in your pub-cache \code{bin} folder. Neither is on PATH by default.
  \end{hint}
\end{frame}

\begin{frame}[fragile,eyebrow=Setup]{Where your code goes}
  \begin{steps}
    \item Branch the team repository:
\begin{shell}
git checkout -b lab-6
\end{shell}
    \item Put the contract in \code{proto/readings/v1/}.
    \item Server code in \code{server/lab6/main.go}. App code in \code{lib/lab6/}.
    \item Run the server once before writing any Dart code.
  \end{steps}
  \begin{hint}[Order of work]
    Contract, then server, then client. Do not write a widget before \code{go run} prints a listening line in your terminal.
  \end{hint}
\end{frame}

\begin{frame}[fragile,eyebrow=Setup]{The contract: readings.proto}
\begin{codeblock}
syntax = "proto3";
package readings.v1;
message Reading {
  string device_id    = 1;
  double value         = 2;
  int64  taken_at_unix = 3;
}
message DeviceRequest { string device_id = 1; }
service Readings {
  rpc GetLatest(DeviceRequest) returns (Reading);
  rpc StreamReadings(DeviceRequest) returns (stream Reading);
}
\end{codeblock}
  \vskip 1mm
  \muted{GetLatest returns one Reading. StreamReadings returns a stream of them.}
\end{frame}

\begin{frame}[fragile,eyebrow=Task I]{Generate the stubs, run the server}
  \begin{columns}[T]
    \begin{column}{0.55\textwidth}
      \begin{steps}
        \item Run protoc for Go, then for Dart (the exact commands are next).
        \item Add the packages: \code{flutter pub add grpc protobuf}.
        \item Start the server: \code{go run ./server/lab6}.
        \item Confirm it prints the port it is listening on.
      \end{steps}
    \end{column}
    \begin{column}{0.41\textwidth}
      \kv{Done when}{}
      \begin{checklist}
        \item \code{gen/readings/v1/} holds two generated Go files.
        \item \code{lib/gen/readings/v1/} holds four Dart files.
        \item The server keeps running with no error printed.
        \item \code{flutter pub get} resolves cleanly.
      \end{checklist}
    \end{column}
  \end{columns}
  \begin{takeaway}{Generated code is build output.}
    Never hand-edit it. Change the .proto file and regenerate instead.
  \end{takeaway}
\end{frame}

\begin{frame}[fragile,eyebrow=Task I]{Generating stubs for Go and Dart}
  \vskip 1mm
  Run protoc once per language, against the same .proto file.
\begin{shell}
protoc -Iproto \
  --go_out=gen      --go_opt=paths=source_relative \
  --go-grpc_out=gen --go-grpc_opt=paths=source_relative \
  proto/readings/v1/readings.proto

protoc -Iproto --dart_out=grpc:lib/gen \
  proto/readings/v1/readings.proto
\end{shell}
  \muted{Commit the generated files, or run protoc in CI. Either way, pick one and write it down.}
\end{frame}

\begin{frame}[fragile,eyebrow=Task I]{A Go server: struct and GetLatest}
  \begin{columns}[T]
    \begin{column}{0.6\textwidth}
\begin{golang}
type server struct {
    pb.UnimplementedReadingsServer
    latest *pb.Reading
}

func (s *server) GetLatest(
    ctx context.Context, req *pb.DeviceRequest,
) (*pb.Reading, error) {
    if s.latest == nil {
        return nil, status.Error(codes.NotFound, "empty")
    }
    return s.latest, nil
}
\end{golang}
    \end{column}
    \begin{column}{0.37\textwidth}
      \begin{itemize}
        \item \textbf{Embed the Unimplemented type.} \muted{It gives every method a stub, so the server compiles before GetLatest exists.}
        \item \textbf{No reading yet is an error,} \muted{not a zero value: NotFound tells the client empty and absent apart.}
      \end{itemize}
    \end{column}
  \end{columns}
\end{frame}

\begin{frame}[fragile,eyebrow=Task I]{A Go server: the stream and main}
  \begin{columns}[T]
    \begin{column}{0.58\textwidth}
\begin{golang}
func (s *server) StreamReadings(
    req *pb.DeviceRequest,
    stream pb.Readings_StreamReadingsServer,
) error {
    ctx := stream.Context()
    for ctx.Err() == nil {
        time.Sleep(time.Second)
        stream.Send(s.latest)
    }
    return nil
}
\end{golang}
    \end{column}
    \begin{column}{0.39\textwidth}
\begin{golang}
func main() {
    lis, _ := net.Listen(
        "tcp", ":50051")
    gs := grpc.NewServer()
    pb.RegisterReadingsServer(
        gs, &server{})
    gs.Serve(lis)
}
\end{golang}
    \end{column}
  \end{columns}
  \muted{StreamReadings loops until the client's context is done. main wires the implementation to a TCP listener.}
\end{frame}

\begin{frame}[fragile,eyebrow=Task II]{A unary call with a deadline}
  \begin{columns}[T]
    \begin{column}{0.55\textwidth}
      \begin{steps}
        \item Add a debug flag that sleeps 3 seconds before \code{GetLatest} returns.
        \item Call \code{getLatest} from Flutter with a 2-second \code{CallOptions} timeout.
        \item Catch \code{GrpcError} and show its status code on screen.
        \item Confirm the call fails while the server is still asleep.
      \end{steps}
    \end{column}
    \begin{column}{0.41\textwidth}
      \kv{Done when}{}
      \begin{checklist}
        \item The screen shows an error, not a stuck spinner.
        \item The error code is \code{deadlineExceeded}.
        \item Turning off the sleep makes the same call succeed.
      \end{checklist}
    \end{column}
  \end{columns}
  \begin{takeaway}{A deadline is not optional.}
    Without one, a slow server holds your UI state forever.
  \end{takeaway}
\end{frame}

\begin{frame}[fragile,eyebrow=Task II]{A Dart client: the unary call}
  \begin{columns}[T]
    \begin{column}{0.6\textwidth}
\begin{dart}
final channel = ClientChannel(
    '10.0.2.2', port: 50051,
    options: const ChannelOptions(
        credentials: ChannelCredentials.insecure()));
final client = ReadingsClient(channel); try {
  final r = await client.getLatest(
    DeviceRequest()..deviceId = id,
    options: CallOptions(
        timeout: const Duration(seconds: 2)));
  setState(() => _value = r.value);
} on GrpcError catch (e) {
  setState(() => _error = e.code); }
\end{dart}
    \end{column}
    \begin{column}{0.37\textwidth}
      \begin{itemize}
        \item \textbf{Insecure means plaintext,} \muted{fine for a lab server on your own machine, never for a real deployment.}
        \item \textbf{10.0.2.2} \muted{is the Android emulator's address for your host machine, not \code{localhost}.}
      \end{itemize}
    \end{column}
  \end{columns}
\end{frame}

\begin{frame}[fragile,eyebrow=Task III]{The streaming screen, with a cancel}
  \begin{columns}[T]
    \begin{column}{0.55\textwidth}
      \begin{steps}
        \item Call \code{streamReadings} with the device id.
        \item Listen on the stream and update the screen on each \code{Reading}.
        \item Add a Cancel button that stops the subscription.
        \item Show a stopped label once canceled.
      \end{steps}
    \end{column}
    \begin{column}{0.41\textwidth}
      \kv{Done when}{}
      \begin{checklist}
        \item The value on screen changes about once a second.
        \item Cancel stops the updates immediately.
        \item Nothing is logged as an error after a clean cancel.
      \end{checklist}
    \end{column}
  \end{columns}
\end{frame}

\begin{frame}[fragile,eyebrow=Task III]{A Dart client: the stream and cancel}
  \begin{columns}[T]
    \begin{column}{0.55\textwidth}
\begin{dart}
final call = client.streamReadings(
  DeviceRequest()..deviceId = id,
);
final sub = call.listen(
  (r) => setState(() => _value = r.value),
  onError: (_) => setState(() => _live = false),
);

void _cancel() {
  call.cancel();   // stop the stream
  sub.cancel();
  setState(() => _live = false);
}
\end{dart}
    \end{column}
    \begin{column}{0.42\textwidth}
      \begin{itemize}
        \item \textbf{call is a ResponseStream.} \muted{Canceling it tells the server to stop sending, not only to stop listening locally.}
        \item \textbf{The timeout still applies.} \muted{Pick a generous one for a stream you plan to keep open for a while.}
      \end{itemize}
    \end{column}
  \end{columns}
\end{frame}

\begin{frame}[eyebrow=Acceptance]{What the grader will do}
  \vskip 1mm
  \begin{columns}[T]
    \begin{column}{0.48\textwidth}
      \kv{The demo}{}
      \begin{checklist}
        \item The stream updates live, about once a second.
        \item Pressing Cancel stops it immediately.
        \item A 2-second deadline fails \code{getLatest} against a server sleeping 3 seconds, with \code{deadlineExceeded}.
      \end{checklist}
    \end{column}
    \begin{column}{0.48\textwidth}
      \kv{In the repository}{}
      \muted{The .proto file, the server, the app code, and a README naming your device id scheme and the port you used.\par\vskip 1.5mm
      One commit per task, on branch \code{lab-6}.}
    \end{column}
  \end{columns}
\end{frame}

\begin{frame}[eyebrow=Troubleshooting]{Errors you will meet, and their fixes}
  \trouble{protoc-gen-go: program not found}{Not on PATH. Add both plugin bin folders to PATH.}
  \trouble{Target of URI doesn't exist: readings.pb.dart}{Match \code{--dart\_out} to where the app imports it.}
  \trouble{Connection refused to localhost:50051}{Emulator localhost is itself. Use \code{10.0.2.2}.}
  \trouble{UNAVAILABLE: failed to connect}{The server isn't running, or the port is wrong.}
  \trouble{CLEARTEXT communication not permitted}{Android blocks it. Set \code{usesCleartextTraffic}.}
\end{frame}

\closingframe{Push before you leave.}{Branch lab-6 · one commit per task · README with device id scheme and port}

\end{document}
